\section{Solução}

Nesta seção, será desinvolvido o exercício proposto de Equações Diferenciais
Ordinárias, com o objetivo de determinar a solução complementar $y_c$ e a
solução particular $y_p$ de uma equação diferencial linear não homogênea de
sexta ordem. Para isso, será empregado o método dos coeficientes
indeterminados, utilizando a abordagem por superposição, que permite decompor o
termo não homogêneo em parcelas e determinar uma solução particular
correspondente a cada uma delas. A solução particular geral será obtida pela
soma das soluções particulares individuais, enquanto a solução complementar
será determinada a partir da equação homogênea associada. Conforme solicitado
no enunciado, não serão calculados os valores das constantes arbitrárias.

\question{
  Encontre $y_c$ e $y_p$ da EDO, utilizando o método dos coeficientes
  indeterminados - abordagem por superposição. Não calcule as constantes.

  \begin{equation}
  \begin{aligned}
    \frac{d^6 y}{dt^6}
    - \frac{d^5 y}{dt^5}
    - 4\frac{d^4 y}{dt^4}
    + 4\frac{d^3 y}{dt^3}
    + 11\frac{d^2 y}{dt^2}
    + 5\frac{dy}{dt}
    &= -e^{-x}
    + x^2 e^{2x}\cos(x)
      + x^3
      - x e^{2x}\sin(x) \\
      &\quad + x^2
      + \cos(x)
      + \cos(2x)
      + x^3 e^{-x}
    + x^4
  \end{aligned}
  \end{equation}
}

\answer{
  A solução geral de uma equação diferencial linear não homogênea é dada pela
  soma da solução complementar $y_c$ com uma solução particular $y_p$, isto é,
  
  \begin{equation}
    \boxed{
      y_g = y_c + y_p.
    }
  \end{equation}
  
  Inicialmente, determinaremos a solução complementar $y_c$, obtida a partir
  da equação diferencial homogênea associada à equação original. Para isso,
  consideramos apenas os termos que envolvem a função incógnita $y$ e suas
  derivadas, igualando a expressão resultante a zero:
  
  \begin{equation*}
    \frac{d^6 y}{dt^6}
    - \frac{d^5 y}{dt^5}
    - 4\frac{d^4 y}{dt^4}
    + 4\frac{d^3 y}{dt^3}
    + 11\frac{d^2 y}{dt^2}
    + 5\frac{dy}{dt}
    = 0.
  \end{equation*}
  
  Em seguida, para determinar a solução complementar, substituímos cada
  derivada de ordem $n$ pelo termo correspondente $m^n$, obtendo a equação
  característica associada:
  
  \begin{equation}
    m^6 - m^5 - 4m^4 + 4m^3 + 11m^2 + 5m = 0.
    \label{eq:carac}
  \end{equation}
  

  Para determinar suas raízes, aplicaremos o dispositivo de Briot--Ruffini
  ao polinômio característico da Equação \ref{eq:carac}. Inicialmente,
  colocando $m$ em evidência, obtemos
  
  \begin{equation*}
    m\left(m^5 - m^4 - 4m^3 + 4m^2 + 11m + 5\right)=0.
  \end{equation*}
  
  Uma das raízes da equação característica é, portanto, $m_1=0$. Para
  determinar as demais raízes, aplicaremos o dispositivo de Briot--Ruffini
  sucessivamente, utilizando $m=-1$ como raiz do polinômio restante. As três
  aplicações são apresentadas a seguir:
  
  \begin{equation*}
  \renewcommand{\arraystretch}{1.3}
  \setlength{\arraycolsep}{5pt}
  \begin{array}{c|rrrrrr}
  -1 & 1 & -1 & -4 & 4 & 11 & 5 \\
     &   & -1 & 2 & 2 & -6 & -5 \\ \hline
     & 1 & -2 & -2 & 6 & 5 & 0
  \end{array}
  \qquad
  \begin{array}{c|rrrrr}
  -1 & 1 & -2 & -2 & 6 & 5 \\
     &   & -1 & 3 & -1 & -5 \\ \hline
     & 1 & -3 & 1 & 5 & 0
  \end{array}
  \qquad
  \begin{array}{c|rrrr}
  -1 & 1 & -3 & 1 & 5 \\
     &   & -1 & 4 & -5 \\ \hline
     & 1 & -4 & 5 & 0
  \end{array}
  \end{equation*}
  
  Após as três aplicações, obtemos o polinômio restante
  $m^2-4m+5$. Para determinar suas raízes, utilizamos a fórmula de
  Bhaskara:
  
  \begin{equation*}
  \begin{aligned}
    m
    &= \frac{-(-4)\pm\sqrt{(-4)^2-4(1)(5)}}{2(1)} \\
    &= \frac{4\pm\sqrt{-4}}{2} \\
    &= 2\pm i.
  \end{aligned}
  \end{equation*}
  
  Portanto, a fatoração do polinômio característico é dada por
  
  \begin{equation*}
    m(m+1)^3(m^2-4m+5)=0,
  \end{equation*}
  
  cujas raízes são
  
  \begin{equation}
    m_1=0,\qquad
    m_2=m_3=m_4=-1,\qquad
    m_5=2+i,\qquad
    m_6=2-i.
  \end{equation}
  
  Com as raízes determinadas, podemos construir a solução complementar
  $y_c$ a partir das contribuições de cada raiz ou conjunto de raízes da
  equação característica.
  
  Inicialmente, para a raiz $m_1=0$, temos a contribuição
  
  \begin{equation*}
    y_{c_1}=C_1e^{0t}=C_1.
  \end{equation*}
  
  Em seguida, como a raiz $m=-1$ possui multiplicidade três, correspondente
  às raízes $m_2=m_3=m_4=-1$, sua contribuição é dada por
  
  \begin{equation*}
    y_{c_2}=(C_2+C_3t+C_4t^2)e^{-t}.
  \end{equation*}
  
  Por fim, para o par de raízes complexas conjugadas
  $m_5=2+i$ e $m_6=2-i$, utilizamos a forma real da solução associada
  a raízes complexas, obtendo
  
  \begin{equation*}
    y_{c_3}=e^{2t}\left(C_5\cos t+C_6\sin t\right).
  \end{equation*}
  
  Assim, a solução complementar é obtida pela soma das três contribuições:
  
  \begin{equation}
    \boxed{
      y_c=C_1+(C_2+C_3t+C_4t^2)e^{-t}
      +e^{2t}\left(C_5\cos t+C_6\sin t\right)
    }
  \end{equation}

  \newpage

  Para determinar a solução particular $y_p$, utilizaremos a Tabela
  \ref{tab:solucoes_particulares}, que apresenta as formas de tentativa
  correspondentes aos diferentes tipos de funções presentes no termo não
  homogêneo $g(x)$. Pela abordagem da superposição, decompomos $g(x)$ em
  parcelas e agrupamos os termos de mesma natureza, construindo uma tentativa
  particular para cada grupo.

  \begin{table}[htbp]
    \centering
    \caption{Tentativas para soluções particulares}
    \label{tab:tentativas-solucoes-particulares}
    \renewcommand{\arraystretch}{1.2}
    \begin{tabular}{|c|l|}
      \hline
      $g(x)$ & $y_p$ \\ \hline
      $1$ (qualquer constante) & $A$ \\ \hline
      $x+5$ & $Ax+B$ \\ \hline
      $3x^2-2$ & $Ax^2+Bx+C$ \\ \hline
      $x^3-x+1$ & $Ax^3+Bx^2+Cx+D$ \\ \hline
      $\sin(4x)$ & $A\cos(4x)+B\sin(4x)$ \\ \hline
      $\cos(4x)$ & $A\cos(4x)+B\sin(4x)$ \\ \hline
      $e^{5x}$ & $Ae^{5x}$ \\ \hline
      $(9x-2)e^{5x}$ & $(Ax+B)e^{5x}$ \\ \hline
      $x^2e^{5x}$ & $(Ax^2+Bx+C)e^{5x}$ \\ \hline
      $e^{3x}\sin(4x)$ & $Ae^{3x}\sin(4x)+Be^{3x}\cos(4x)$ \\ \hline
      $5x^2\sin(4x)$ & $(Ax^2+Bx+C)\cos(4x)+(Dx^2+Ex+F)\sin(4x)$ \\ \hline
      $xe^{3x}\cos(4x)$ & $(Ax+B)e^{3x}\cos(4x)+(Cx+D)e^{3x}\sin(4x)$ \\ \hline
    \end{tabular}
    \label{tab:solucoes_particulares}
  \end{table}
  
  O termo não homogêneo da equação é
  
  \begin{equation}
    \begin{aligned}
      g(x) ={}&
      -e^{-x}
      + x^2 e^{2x}\cos(x)
      - x e^{2x}\sin(x) \\
      &+ x^3 + x^2 + x^4
      + \cos(x) + \cos(2x)
      + x^3e^{-x}.
    \end{aligned}
  \end{equation}
  
  Agrupando os termos semelhantes, obtemos as seguintes parcelas:
  
  \begin{align*}
    g_1(x) &= -e^{-x}, \\
    g_2(x) &= x^2e^{2x}\cos(x)-xe^{2x}\sin(x), \\
    g_3(x) &= x^4+x^3+x^2, \\
    g_4(x) &= \cos(x)+\cos(2x), \\
    g_5(x) &= x^3e^{-x}.
  \end{align*}
  
  Assim, pela propriedade da superposição, a solução particular pode ser
  escrita como:
  
  \begin{equation*}
    y_p = y_{p_1}+y_{p_2}+y_{p_3}+y_{p_4}+y_{p_5}.
  \end{equation*}
  
  Para cada parcela, propomos as seguintes formas de tentativa:
  
  \begin{enumerate}
    \item Para $g_1(x)=-e^{-x}$, identificamos uma função exponencial da
    forma $e^{mx}$, com $m=-1$. Ao consultar as raízes da equação
    característica, verificamos que $m=-1$ é uma raiz de multiplicidade
    três, cuja contribuição para a solução complementar é
    $(C_1+C_2x+C_3x^2)e^{-x}$. Como a tentativa inicial $Ae^{-x}$
    já está contida na solução complementar, devemos multiplicá-la por
    $x^3$, correspondente à multiplicidade da raiz. Assim, propomos
  
    \begin{equation*}
      y_{p_1}=Ax^3e^{-x}.
    \end{equation*}
  
    \item Para
    $g_2(x)=x^2e^{2x}\cos(x)-xe^{2x}\sin(x)$, identificamos termos
    exponenciais associados a funções trigonométricas. A expressão
    $e^{2x}\cos(x)$ e $e^{2x}\sin(x)$ corresponde ao par de raízes
    complexas $m=2\pm i$. Como essas raízes são simples, mas já aparecem
    na solução complementar por meio de
    $e^{2x}(C_4\cos x+C_5\sin x)$, a tentativa particular deve ser
    multiplicada por $x$. Além disso, como os coeficientes polinomiais
    possuem grau máximo dois, utilizamos polinômios de grau dois tanto
    para o seno quanto para o cosseno. Portanto,
  
    \begin{equation*}
      \begin{aligned}
        y_{p_2}={}&
        x e^{2x}\left[
        (Ax^2+Bx+C)\cos(x)\right.\\
        &\left.\qquad +(Dx^2+Ex+F)\sin(x)
        \right].
      \end{aligned}
    \end{equation*}
  
    \item Para $g_3(x)=x^4+x^3+x^2$, temos um polinômio de grau quatro.
    Como o termo não homogêneo corresponde à exponencial $e^{0x}=1$,
    devemos verificar se $m=0$ é raiz da equação característica.
    De fato, $m=0$ é uma raiz simples e já está representada pela
    constante presente em $y_c$. Portanto, a tentativa polinomial
    usual deve ser multiplicada por $x$. Assim, propomos
  
    \begin{equation*}
      y_{p_3}=x(Ax^4+Bx^3+Cx^2+Dx+E).
    \end{equation*}
  
    \item Para $g_4(x)=\cos(x)+\cos(2x)$, identificamos duas frequências
    trigonométricas distintas. A função $\cos(x)$ está associada às
    raízes complexas $m=\pm i$, enquanto $\cos(2x)$ está associada
    às raízes $m=\pm 2i$. Como nenhuma dessas raízes pertence ao
    conjunto de raízes da equação característica, não há
    ressonância e, portanto, não precisamos multiplicar as tentativas
    por $x$. Dessa forma, propomos
  
    \begin{equation*}
      \begin{aligned}
        y_{p_4}={}&A\cos(x)+B\sin(x)\\
        &+C\cos(2x)+D\sin(2x).
      \end{aligned}
    \end{equation*}
  
    \item Para $g_5(x)=x^3e^{-x}$, temos novamente uma função
    exponencial associada a $m=-1$, acompanhada de um polinômio de
    grau três. Como essa raiz possui multiplicidade três na equação
    característica, a tentativa usual
    $e^{-x}(Ax^3+Bx^2+Cx+D)$ apresenta termos que se sobrepõem à
    solução complementar. Para eliminar essa sobreposição, devemos
    multiplicar todo o polinômio por $x^3$. Logo,

    \begin{equation*}
      y_{p_5}=x^3e^{-x}(Ax^3+Bx^2+Cx+D).
    \end{equation*}
  \end{enumerate}
  
  Somando as tentativas individuais, obtemos a forma geral de $y_p$:

  \begin{equation*}
    \begin{aligned}
      y_p ={}&
      Ax^3e^{-x} \\[4pt]
      &+ xe^{2x}\left[
        (Bx^2+Cx+D)\cos(x)
        +(Ex^2+Fx+G)\sin(x)
      \right] \\[4pt]
      &+ x(Hx^4+Ix^3+Jx^2+Kx+L) \\[4pt]
      &+ M\cos(x)+N\sin(x)
        +O\cos(2x)+P\sin(2x) \\[4pt]
      &+ x^3e^{-x}(Qx^3+Rx^2+Sx+T).
    \end{aligned}
  \end{equation*}

  Ao somar as tentativas particulares individuais, observamos que os termos
  associados a $e^{-x}$, provenientes de $g_1(x)$ e $g_5(x)$, pertencem à mesma
  família de soluções. Portanto, podemos agrupá-los em uma única tentativa
  particular, evitando a redundância de coeficientes. Dessa forma,
  simplificamos a expressão de $y_p$, mantendo todas as contribuições
  necessárias:

  \begin{equation}
    \boxed{
      \begin{aligned}
        y_p ={}&
        x^3e^{-x}(Ax^3+Bx^2+Cx+D) \\[4pt]
        &+ xe^{2x}\left[
          (Ex^2+Fx+G)\cos(x)
          +(Hx^2+Ix+J)\sin(x)
        \right] \\[4pt]
        &+ x(Kx^4+Lx^3+Mx^2+Nx+O) \\[4pt]
        &+ P\cos(x)+Q\sin(x)
          +R\cos(2x)+S\sin(2x).
      \end{aligned}
    }
  \end{equation}

  \newpage

  Já com $y_c$ e $y_p$, podemos obter a solução geral $y_g$, uma vez que o
  exercício não nos pede para calcular as constantes arbitrárias. Assim, basta
  somar a solução complementar à solução particular, obtendo $y_g = y_c + y_p$:
  
  \begin{equation}
    \boxed{
      \begin{aligned}
        y_g ={}&
        C_1+(C_2+C_3t+C_4t^2)e^{-t}
        +e^{2t}\left(C_5\cos t+C_6\sin t\right) \\[6pt]
        &+x^3e^{-x}(Ax^3+Bx^2+Cx+D) \\[4pt]
        &+xe^{2x}\left[
          (Ex^2+Fx+G)\cos x
          +(Hx^2+Ix+J)\sin x
        \right] \\[4pt]
        &+x(Kx^4+Lx^3+Mx^2+Nx+O) \\[4pt]
        &+P\cos x+Q\sin x
          +R\cos(2x)+S\sin(2x).
      \end{aligned}
    }
  \end{equation}

}
